% 小行星（综述）
% license CCBYSA3
% type Wiki

本文根据 CC-BY-SA 协议转载翻译自维基百科\href{https://en.wikipedia.org/wiki/Asteroid}{相关文章}。

\begin{figure}[ht]
\centering
\includegraphics[width=8cm]{./figures/9fdf35f418301f87.png}
\caption{被探测访问过的小行星影像展示其差异：（上排）433 Eros 与 243 Ida 及其卫星 Dactyl，（下排）Ceres 与 101955 Bennu。尺寸不按比例绘制。} \label{fig_Astero_1}
\end{figure}
小行星是一类小行星体（minor planet），其尺寸大于流星体（meteoroid），但既不是行星，也不是已确认为彗星的天体。它们位于太阳系内部轨道运行，或与木星共享轨道（如特洛伊小行星）。小行星由岩石、金属或冰质物质构成，不具备大气层，并可大致分类为 C 型（碳质）、M 型（金属）、S 型（硅酸盐）。小行星的大小与形状差异极大，从直径不足一千米的松散碎石堆，到直径接近 1000 千米、被归类为矮行星的 Ceres。若天体在受到太阳辐射加热时产生彗发（尾），则被归类为彗星而非小行星，尽管近期观测显示这两类天体之间可能存在连续光谱。\(^\text{[1][2]}\)

在约一百万颗已知小行星中，\(^\text{[3]}\)数量最多的位于火星与木星轨道之间、距离太阳约 2 至 4 AU 的区域，称为主小行星带（main asteroid belt）。所有小行星总质量仅约为地球月球质量的 3\%。主带中大多数小行星沿轻微椭圆、稳定的轨道运行，方向与地球一致，绕太阳一周约需 3 至 6 年。\(^\text{[4]}\)

小行星最早通过地面观测方式进行研究，首次近距离观测来自伽利略号（Galileo）探测器。其后，美国 NASA 与日本 JAXA 相继发射了多次专门的小行星探测任务，并持续规划新的任务。NASA 的 NEAR Shoemaker 研究了 Eros，Dawn 探测了 Vesta 与 Ceres；JAXA 的隼鸟号（Hayabusa）与隼鸟 2 号分别研究并回收了 Itokawa 与 Ryugu 的样本；OSIRIS{-}REx 探测 Bennu，并于 2020 年采样，2023 年将样本送回地球。NASA 的 Lucy（2021 年发射）将研究十颗不同小行星，其中两颗来自主带、八颗为木星特洛伊小行星。Psyche 探测器于 2023 年 10 月发射，目标是金属小行星 Psyche。欧洲航天局的 Hera 于 2024 年 10 月发射，旨在研究 DART 碰撞后的结果。中国国家航天局的天问二号于 2025 年 5 月发射，\(^\text{[5]}\)其任务为探测共轨近地小行星 469219 Kamoʻoalewa 与活跃小行星 311P/PanSTARRS，并对 Kamoʻoalewa 的风化层（regolith）进行采样。\(^\text{[6]}\)

近地小行星若与地球发生撞击，可能造成灾难性影响，其中最著名的案例是 Chicxulub 撞击事件，被普遍认为导致了白垩纪–古近纪生物大灭绝。为应对这一潜在威胁，作为实验性防御措施，2022 年 9 月，双小行星偏转测试（Double Asteroid Redirection Test, DART）探测器通过撞击成功改变了无威胁小行星 Dimorphos 的轨道。
\subsection{术语}
\begin{figure}[ht]
\centering
\includegraphics[width=8cm]{./figures/795a52138582bb97.png}
\caption{一幅按相同尺度制作的组合图像，展示了 2012 年前以高分辨率成像的小行星。从大到小依次为：4 Vesta、21 Lutetia、253 Mathilde、243 Ida 及其卫星 Dactyl、433 Eros、951 Gaspra、2867 Šteins、25143 Itokawa。} \label{fig_Astero_2}
\end{figure}
\begin{figure}[ht]
\centering
\includegraphics[width=8cm]{./figures/1f1fc84d80338171.png}
\caption{} \label{fig_Astero_3}
\end{figure}
2006 年，国际天文学联合会（IAU）引入了当前更为推荐的宽泛术语“太阳系小天体”（small Solar System body），将其定义为太阳系中既不是行星、也不是矮行星、亦不是天然卫星的天体；其中包括小行星、彗星以及更近期发现的其他类型天体[7]。根据 IAU 表述，“术语『小行星（minor planet）』仍可继续使用，但一般情况下将优先采用『太阳系小天体（Small Solar System Body）』”[8]。

在历史上，第一颗被发现的小行星——Ceres——最初曾被视为一颗新行星[a]。随后，天文学家又陆续发现了其他相似天体。在当时的观测设备条件下，它们呈现为类似恒星的光点，几乎看不到行星盘结构，但由于具有明显的视运动，又可与恒星区分开来。这促使天文学家威廉·赫歇尔（Sir William Herschel）提出了“asteroid”一词[b]，该词源自希腊语 ἀστεροειδής（asteroeidēs），意为“类星”“似星状”，并由古希腊语 ἀστήρ（astēr，“星、行星”）派生。在 19 世纪后半叶初期，“asteroid”与“planet”（有时并未特别限定为“minor”）仍被交替使用[c]。

传统上，围绕太阳运行的小型天体被划分为彗星、小行星或流星体，其中直径小于一米的天体被称为流星体。术语“asteroid”从未被正式定义[13]，但在非正式使用中可指“围绕太阳运行、形状不规则、呈岩质组成且不符合 IAU 定义中行星或矮行星条件的天体”[14]。小行星与彗星的主要区别在于，彗星因近表层冰在太阳辐射作用下发生升华而形成彗发（尾）。少数天体最初被归类为小行星，但后续观测显示存在彗状活动。相反，一些（或许全部）彗星最终会耗尽其表面挥发性冰，从而演化为类小行星天体。另一区分方式是轨道偏心率：彗星通常具有比大多数小行星更高的轨道偏心率；而高度偏心的小行星极可能为休眠或灭绝的彗星[15]。
```
